\documentclass[11pt]{article}
\usepackage[margin=25mm]{geometry}
\usepackage{fontspec,unicode-math}
\setmainfont{Latin Modern Roman}
\setmathfont{Latin Modern Math}
\usepackage{amsmath,amssymb,mathtools}
\usepackage{xurl,hyperref}
\hypersetup{colorlinks=true,urlcolor=blue,linkcolor=blue}
\setcounter{secnumdepth}{0}
\setlength{\parindent}{0pt}
\setlength{\parskip}{5pt}
\setlength{\emergencystretch}{3em}
\newcommand{\sref}[2]{\hyperlink{source:#1:#2}{[S#2]}}
\newcommand{\cataloguescope}[1]{\paragraph{Source question.}#1}
\newcommand{\floor}[1]{\left\lfloor#1\right\rfloor}
\begin{document}
\section{Real Tau Conjecture for Depth-Four Circuits}\label{statement}

\subsection{Definitions and mathematical statement}

\paragraph{Shared notation.}Standard number systems are denoted by $\mathbb N,\mathbb Z,\mathbb Q,\mathbb R,\mathbb C$; index ranges, computational models and specialized notation follow the cited problem source and the definitions below.


A sparse univariate polynomial has few nonzero monomials, with exponents encoded as prescribed. The stated depth-four form sums products of such sparse polynomials; real roots are counted without multiplicity.

An arithmetic circuit is a directed acyclic graph whose input nodes are variables or field constants and whose internal nodes perform additions and multiplications (with subtraction allowed when stated). Size counts gates; depth is the maximum number of gates on an input-output path.

A Boolean function maps $\{0,1\}^n$ to $\{0,1\}$. A Boolean circuit is an acyclic network of logical gates; a formula is a circuit with no sharing of intermediate gate outputs. Size and depth measure its number of gates and longest gate path.
\paragraph{Statement recorded in the supplied source.}
Is the number of distinct real zeros of every univariate real polynomial represented by a depth-four arithmetic circuit, consisting of a sum of products of sparse polynomials, bounded by a polynomial in the circuit size? \sref{statement}{2}
\paragraph{Scope qualification.}
The size convention must account for the source\textquotesingle s sparsity, products, and exponent encoding; no degree bound is inserted.
\paragraph{Source status.}
The supplied catalogue labels this record partially solved. Its source-status metadata was checked against the supplied snapshot on 2026-10-11. This import does not independently establish current open status or mathematical correctness. The catalogue literature assessment is dated 2026-08-21.
\subsection{Short English statement}
Bound the number of distinct real roots by a polynomial in the displayed arithmetic representation size.
\subsection{Sources}
\begin{itemize}
\item[\hypertarget{source:statement:1}{S1}] Ulam AI and the UnsolvedMath Contributors, UnsolvedMath: A Curated Collection of Open Mathematics Problems, catalogue entry. Catalogue curation and metadata: CC BY 4.0. The supplied source question is retained; definitions and the English summary are editorial additions.
\url{https://huggingface.co/datasets/ulamai/UnsolvedMath}.
\item[\hypertarget{source:statement:2}{S2}] Complexity Theory (2019), Oberwolfach Reports, DOI 10.4171/owr/2018/51, pp. 3025--3080
\url{https://doi.org/10.4171/owr/2018/51}.
\item[\hypertarget{source:statement:3}{S3}] The Real Tau-Conjecture Is True on Average
\url{https://arxiv.org/abs/1806.00417}.
\item[\hypertarget{source:statement:4}{S4}] Oberwolfach Report 51/2018
\url{https://ems.press/content/serial-article-files/46773}.
\textit{Underlying Oberwolfach Reports material retains its source-specific licence, including CC BY-SA 4.0 where applicable; see THIRD\_PARTY\_NOTICES.md.}
\end{itemize}
\paragraph{Preparation and review.}
Original question attributed to the cited source and UnsolvedMath contributors. Codex (GPT-6; exact serving revision not exposed) assisted with definitions, English summaries, source inspection and formatting. No new human mathematical authorship or independent proof verification is asserted. Catalogue identity checked against the complete supplied registry and existing repository statements; known flagged repetitions and obvious identical conjectures excluded. Full mathematical-equivalence/literature review has not been performed. Statement preparation and final review followed contribution prompts C and D; automated checks verify structure and identity, not open status.
\label{end:statement}
\end{document}
